\subsection{Quadratic Forms and Polar Forms}

Throughout this subsection, $\mathcal{V}$ is a vector space over an arbitrary field $\mathbb{F}$.

\begin{definition} \label{def:quadratic-form-and-polar-form} \hlt{Quadratic Forms and Polar Forms}\\
A \hlt{quadratic form} on $\mathcal{V}$ is a map $Q:\mathcal{V}\to\mathbb{F}$ satisfying:
\begin{enumerate}[label=(\roman*)]
\setlength{\itemsep}{0pt}
\item $(\forall\,a\in\mathbb{F})(\forall\,\mathbf{v}\in\mathcal{V})\,Q(a\mathbf{v})=a^2Q(\mathbf{v})$;
\item the \hlt{polar form} $B_Q:\mathcal{V}\times\mathcal{V}\to\mathbb{F}$ defined by $B_Q(\mathbf{u},\mathbf{v})\defeq Q(\mathbf{u}+\mathbf{v})-Q(\mathbf{u})-Q(\mathbf{v})$ is bilinear.
\end{enumerate}
\end{definition}

\begin{proposition} \label{prop:quadratic-forms-symmetric-bilinear-forms} \hlt{Quadratic Forms and Symmetric Bilinear Forms}
\begin{enumerate}[label=(\roman*)]
\setlength{\itemsep}{0pt}
\item If $Q:\mathcal{V}\to\mathbb{F}$ is quadratic, then $\operatorname{Sym}(B_Q)$.
\item Let $\operatorname{char}(\mathbb{F})\neq2$ and let $B:\mathcal{V}\times\mathcal{V}\to\mathbb{F}$ be bilinear with $\operatorname{Sym}(B)$. Define $Q_B(\mathbf{x})\defeq\frac{1}{2}B(\mathbf{x},\mathbf{x})$; $Q_B$ is quadratic and $B_{Q_B}=B$.
\item Let $\operatorname{char}(\mathbb{F})\neq2$ and let $Q:\mathcal{V}\to\mathbb{F}$ be quadratic. Define $Q_{B_Q}(\mathbf{x})\defeq\frac{1}{2}B_Q(\mathbf{x},\mathbf{x})$. Then $Q_{B_Q}=Q$.
\end{enumerate}
\begin{proof}
\begin{enumerate}[label=(\roman*)]
\setlength{\itemsep}{0pt}
\item[$(i)$] $(\forall\,\mathbf{u},\mathbf{v}\in\mathcal{V})\,B_Q(\mathbf{u},\mathbf{v})=B_Q(\mathbf{v},\mathbf{u})$ by commutativity of vector addition. $\operatorname{Sym}(B_Q)$ by this identity.
\item[$(ii)$] $(\forall\,a\in\mathbb{F})(\forall\,\mathbf{x}\in\mathcal{V})\,Q_B(a\mathbf{x})=a^2Q_B(\mathbf{x})$ by bilinearity.\\
$(\forall\,\mathbf{u},\mathbf{v}\in\mathcal{V})\,B_{Q_B}(\mathbf{u},\mathbf{v})=\frac{1}{2}\bigl(B(\mathbf{u},\mathbf{v})+B(\mathbf{v},\mathbf{u})\bigr)=B(\mathbf{u},\mathbf{v})$ by bilinearity and symmetry.\\
$B_{Q_B}=B$ by this identity, and $Q_B$ is quadratic by the homogeneity identity and bilinearity of $B_{Q_B}$.
\item[$(iii)$] $(\forall\,\mathbf{x}\in\mathcal{V})\,Q_{B_Q}(\mathbf{x})=\frac{1}{2}B_Q(\mathbf{x},\mathbf{x})=Q(\mathbf{x})$ by quadratic homogeneity. $Q_{B_Q}=Q$ by this pointwise identity.
\end{enumerate}
\hlt{Structural Commentary}: Symmetry places $B_Q$ among symmetric forms; items (ii)--(iii) establish inverse reconstruction.\\
\hlt{Key Insight}: The identity $B_Q(\mathbf{x},\mathbf{x})=2Q(\mathbf{x})$ makes diagonal evaluation invert polarization.
\end{proof}
\end{proposition}

\begin{proposition} \label{prop:polar-forms-alternating-in-characteristic-two} \hlt{Polar Forms in Characteristic Two}\\
Let $Q:\mathcal{V}\to\mathbb{F}$ be quadratic. If $\operatorname{char}(\mathbb{F})=2$, then $\operatorname{Alt}(B_Q)$.
\begin{proof}
$Q(\mathbf{0})=0$ by quadratic homogeneity at scalar $0$.\\
$(\forall\,\mathbf{x}\in\mathcal{V})\,B_Q(\mathbf{x},\mathbf{x})=Q(2\mathbf{x})-2Q(\mathbf{x})=Q(\mathbf{0})=0$ by $\operatorname{char}(\mathbb{F})=2$. $\operatorname{Alt}(B_Q)$ by the preceding identity.\\
\hlt{Structural Commentary}: Homogeneity supplies $Q(\mathbf{0})=0$; characteristic $2$ then collapses the polar diagonal.\\
\hlt{Key Insight}: The identities $2\mathbf{x}=\mathbf{0}$ and $2Q(\mathbf{x})=0$ force diagonal vanishing.
\end{proof}
\end{proposition}

\begin{proposition} \label{prop:characteristic-two-square-perturbations} \hlt{Square Perturbations in Characteristic Two}\\
Assume $\operatorname{char}(\mathbb{F})=2$.
\begin{enumerate}[label=(\roman*)]
\setlength{\itemsep}{0pt}
\item Let $Q:\mathcal{V}\to\mathbb{F}$ be quadratic and $L:\mathcal{V}\to\mathbb{F}$ linear, and define $L^2(\mathbf{v})\defeq L(\mathbf{v})^2$. Then $Q+L^2$ is quadratic and $B_{Q+L^2}=B_Q$.
\item On $\mathbb{F}^2$, define $Q_0(x,y)\defeq xy$ and $Q_1(x,y)\defeq xy+x^2$.\\
Then $Q_0,Q_1$ are quadratic, $Q_0\neq Q_1$, and $B_{Q_0}=B_{Q_1}=B$ for $B((x,y),(u,v))\defeq xv+uy$.
\end{enumerate}
\begin{proof}
\begin{enumerate}[label=(\roman*)]
\setlength{\itemsep}{0pt}
\item[$(i)$] $(\forall\,a\in\mathbb{F})(\forall\,\mathbf{v}\in\mathcal{V})\,(Q+L^2)(a\mathbf{v})=a^2(Q+L^2)(\mathbf{v})$ by homogeneity of $Q$ and linearity of $L$. $(\forall\,\mathbf{u},\mathbf{v}\in\mathcal{V})\,B_{Q+L^2}(\mathbf{u},\mathbf{v})-B_Q(\mathbf{u},\mathbf{v})=2L(\mathbf{u})L(\mathbf{v})=0$ by linearity of $L$ and characteristic $2$.\\
$B_{Q+L^2}=B_Q$ by this identity, and $Q+L^2$ is quadratic by the homogeneity identity and bilinearity of $B_Q$.
\item[$(ii)$] Define $L:\mathbb{F}^2\to\mathbb{F}$ by $L(x,y)\defeq x$.\\
$(\forall\,a,x,y\in\mathbb{F})\,Q_0(ax,ay)=a^2Q_0(x,y)$ and $B_{Q_0}=B$ by direct expansion.\\
$Q_0$ is quadratic by this homogeneity identity and bilinearity of $B$.\\
$Q_1$ is quadratic and $B_{Q_1}=B_{Q_0}=B$ by $Q_1=Q_0+L^2$ and (i). $Q_0\neq Q_1$ by $Q_0(1,0)=0\neq1=Q_1(1,0)$.
\end{enumerate}
\hlt{Structural Commentary}: Item (i) isolates the polarization kernel; item (ii) realizes it through the coordinate square $x^2$.\\
\hlt{Key Insight}: In characteristic $2$, linear-functional squares are invisible to polarization.
\end{proof}
\end{proposition}

\begin{example} \label{ex:characteristic-two-alternating-versus-nonalternating-symmetric} \hlt{Alternating versus Nonalternating Symmetric Forms}\\
Assume $\operatorname{char}(\mathbb{F})=2$ and let $H=\left[\begin{smallmatrix}0&1\\1&0\end{smallmatrix}\right]$. The associated symmetric bilinear form on $\mathbb{F}^2$ is $B_H((x,y),(u,v))=xv+yu$.\\
The form is alternating by $(\forall\,x,y\in\mathbb{F})\,B_H((x,y),(x,y))=2xy=0$. The form is nonsingular by $H^2=I_2$ and Proposition \ref{prop:square-bilinear-pairing-equivalences}.\\
A diagonal matrix $D$ congruent to $H$ would represent $B_H$ in an ordered basis by Proposition \ref{prop:metric-form-matrix-recovery-and-congruence}(iii).\\
Alternation would force $D=0$. Congruence would instead force $\operatorname{rank}(D)=\operatorname{rank}(H)=2$ by Proposition \ref{prop:metric-form-matrix-recovery-and-congruence}(iv).\\
$H$ is not congruent to any diagonal matrix by this rank contradiction.\\
By contrast, $[1]$ is symmetric and nonalternating because its basis vector has self-pairing $1$.\\
\hlt{Structural Commentary}: The forms represented by $H$ and $[1]$ separate alternating from nonalternating symmetry.\\
\hlt{Key Insight}: Alternation forces a zero diagonal, but nonsingularity forbids the zero matrix.
\end{example}
